\section{Mathematical and Stability Analysis}
Time-discretisation accuracy has two complementary aspects. Consistency describes how the local truncation error decreases with $h$, while stability determines whether local perturbations remain controlled as they propagate through the integration. For the stiff Robertson system, both must be examined: a high-order explicit method may have a small local truncation error while still being restricted by its bounded absolute-stability region.

\subsection{Full and reduced Jacobians}
Differentiating the full Robertson vector field gives
\[
J(y)=\begin{pmatrix}
-0.04&10^4y_3&10^4y_2\\
0.04&-10^4y_3-6\times10^7y_2&-10^4y_2\\
0&6\times10^7y_2&0
\end{pmatrix}.
\]
For example, differentiating $-3\times10^7y_2^2$ gives $-6\times10^7y_2$. Since $y_3=1-y_1-y_2$, direct substitution into the first two equations gives
\[
J_{\rm red}=\begin{pmatrix}
-0.04-10^4y_2&10^4(1-y_1-2y_2)\\[2mm]
0.04+10^4y_2&10^4(y_1+2y_2-1)-6\times10^7y_2
\end{pmatrix}.
\]
For example,
\[
\frac{\partial}{\partial y_2}
\left[10^4y_2(1-y_1-y_2)\right]
=10^4(1-y_1-2y_2),
\]
which fixes the sign structure of the second column.
Conservation implies $(1,1,1)^TJ=0$, a \emph{left} null-vector relation; it does not make $(1,1,1)^T$ a right eigenvector. The full Jacobian is singular and has a structural zero eigenvalue. The two eigenvalues of $J_{\rm red}$ coincide with the two nonzero modes tangent to the invariant plane.

\begin{figure}[H]
\centering\includegraphics[width=\textwidth]{F5_jacobian_spectrum.pdf}
\caption{Jacobian spectrum along the reference trajectory. The two dynamically relevant modes develop widely separated decay rates, while the full Jacobian has a structural zero eigenvalue. A narrow early-time interval contains a complex-conjugate pair with small imaginary parts. The dense early-time panel reveals this feature, which may be missed by the fixed logarithmic output grid; the later separation of real parts supplies the local time-scale evidence for stiffness.}
\label{fig:spectrum}
\end{figure}

\subsection{Stiffness ratio}
Using only the two nonzero eigenvalues,
\[
S(t)=\frac{\max_{\lambda_j\ne0}|\operatorname{Re}\lambda_j(t)|}
{\min_{\lambda_j\ne0}|\operatorname{Re}\lambda_j(t)|}.
\]
Including the structural zero would make $S$ undefined. The exactly degenerate initial state is not interpreted; the plot begins at the first positive logarithmic time. Eigenvalues are evaluated along the same reference trajectory used for error measurement.

\begin{figure}[H]
\centering\includegraphics[width=\textwidth]{F6_stiffness_ratio.pdf}
\caption{Stiffness ratio on logarithmic time. Check values are approximately $3.0\times10^3$ at $10^{-4}$, $5.4\times10^3$ at $10^{-2}$, and $1.58\times10^5$ at $40$. Its growth shows that the separation of decay time scales persists and strengthens after the visible early transient; it is evaluated only from the two nonzero modes.}
\label{fig:ratio}
\end{figure}

\subsection{Absolute stability and frozen bounds}
\label{sec:frozen-bounds}
The absolute-stability region is $\{z:|R(z)|\le1\}$. For explicit Euler this is the disk $|1+z|\le1$; RK4 has a bounded quartic region; implicit Euler includes the left half-plane because $|1/(1-z)|\le1$ there. These statements concern the linear test equation. Along the nonlinear Robertson trajectory, the eigenvalues of a frozen Jacobian provide local values of $z=h\lambda(t)$ and therefore a diagnostic of where stability restrictions may arise.

\begin{figure}[H]
\centering\includegraphics[width=\textwidth]{F4_stability_regions.pdf}
\caption{Absolute-stability regions with the Robertson spectrum $h\lambda(t)$ overlaid. The bounded EE and RK4 regions expose a local step-size restriction from the fast mode, whereas implicit Euler contains the left half-plane. The overlay provides a local stability scale. Accuracy and non-negativity are evaluated separately in the numerical experiments.}
\label{fig:regions}
\end{figure}

At $t=40$, $|\lambda_{\rm fast}|=3392.79$, hence
\[
h_{EE}\lesssim\frac2{3392.79}=5.895\times10^{-4},\qquad
h_{RK4}\lesssim\frac{2.7853}{3392.79}=8.209\times10^{-4}.
\]
The frozen bounds estimate the explicit step-size scale along the trajectory. The empirical sweep in Section~5.2 tests this estimate against the nonlinear computation. The corresponding value $40/h$ is a step-count estimate under the frozen-Jacobian condition.

The experiments evaluate absolute stability, numerical accuracy, and physical admissibility separately. Absolute stability concerns the linear test equation or a locally frozen linearisation, numerical accuracy concerns error relative to the reference trajectory, and physical admissibility includes non-negative concentrations. A method may remain stable while producing an inaccurate or negative solution.

\subsection{Four diagnostics}
Define
\begin{align*}
E_{mass}&=\max_{t\in\mathcal T}|y_1+y_2+y_3-1|,\\
y_{min}&=\min_{i,t\in\mathcal T}y_i(t),\\
E_{global}&=\max_{t\in\mathcal T}\|y_h(t)-y_{ref}(t)\|_\infty,\\
E_{end}&=|y_{2,h}(40)-y_{2,ref}(40)|.
\end{align*}
The four diagnostics describe different properties of the solution. Since $40\in\mathcal T$,
\[
E_{end}\le E_{global}.
\]
Thus a full-state pass implies an endpoint pass under the same absolute threshold; the converse is false. Conservation does not imply positivity or accuracy.

\begin{figure}[H]
\centering\includegraphics[width=\textwidth]{F9_four_track_diagnostics.pdf}
\caption{Four diagnostic tracks. Conservation alone cannot certify accuracy: the invariant remains useful for detecting implementation errors, although its near-exact preservation does not establish time-discretisation accuracy. The explicit-Euler counterexample at $h=0.003$ has conservation defect about $1.1\times10^{-13}$ while the minimum concentration is about $-921.8$ and the error about $1.08\times10^{10}$.}
\label{fig:diagnostics}
\end{figure}
The $2.1\times10^{-12}$ line records the difference between numerical reference solvers. It measures cross-solver agreement rather than the unknown true solution error, so differences below this level are reported without assigning additional significant digits.
